% ============================================================
%  ch2_I_a —— 第一章 证明的原理（原书 pp. 3--9，即 PDF 第 4--10 页）
%  底本：L. Lafforgue, Cours à l'Institut Tata sur les chtoucas de
%  Drinfeld et la correspondance de Langlands, IHES/M/02/45 (2002)
% ============================================================

\section{证明的原理}

% ------------------------------------------------------------
%  PAGE 4（原书 p. 3）
% ------------------------------------------------------------
\origpage{3}以下始终考虑：

$X$ = 光滑（几何连通）射影曲线 $/\mathbb{F}_q$

$F = F(X)$ = $X$ 上的有理函数域

$|X| = \{\text{$X$ 的闭点}\}$ \qquad $|X| \ni x \mapsto F_x \supset O_x$

$\mathbb{A} = \prod_{x \in |X|} F_x \supset O_{\mathbb{A}} = \prod_{x \in |X|} O_x$

$\deg : \mathbb{A}^\times \to \mathbb{Z}$ \qquad
$\mathbb{A}^{\times 0} = \mathrm{Ker}[\mathbb{A}^\times \xrightarrow{\deg} \mathbb{Z}]$

\[ \mathrm{GL}_r(\mathbb{A})^0 = \mathrm{Ker}[\mathrm{GL}_r(\mathbb{A}) \xrightarrow{\det} \mathbb{A}^\times \xrightarrow{\deg} \mathbb{Z}]. \]

取 $a \in \mathbb{A}^\times$，$\deg(a) \neq 0$。

$F^\times\backslash \mathbb{A}^{\times 0}$、$F^\times\backslash \mathbb{A}^\times/a^{\mathbb{Z}}$ 都是紧的。

当 $r \ge 2$ 时，$\mathrm{GL}_r(F)\backslash \mathrm{GL}_r(\mathbb{A})^0$、
$\mathrm{GL}_r(F)\backslash \mathrm{GL}_r(\mathbb{A})/a^{\mathbb{Z}}$ 具有有限体积，但不紧。

$\{\pi\}_r = \{\text{$\mathrm{GL}_r(\mathbb{A})$ 的尖点自守表示，其中心特征 $\chi_\pi$ 为有限阶}\}$

$\{\pi\}_r \ni \pi$

$|X| \ni x$ 非分歧 $\mapsto z_1(\pi_x), \ldots, z_r(\pi_x)$：Hecke 特征值。

$\{\sigma\}_r = \{\text{$G_F$ 的 $r$ 维不可约 $\ell$-进表示，其行列式为有限阶}\}$

$\{\sigma\}_r \ni \sigma$

$|X| \ni x$ 非分歧 $\mapsto z_1(\sigma_x), \ldots, z_r(\sigma_x)$：$\mathrm{Frob}_x$ 的特征值。

\begin{theorem*}
(i)$_r$\ 对任意 $\pi \in \{\pi\}_r$，存在唯一的 $\sigma = \sigma_\pi \in \{\sigma\}_r$，它在 $\pi$ 非分歧之处非分歧，且
\[ \{z_1(\pi_x), \ldots, z_r(\pi_x)\} = \{z_1(\sigma_x), \ldots, z_r(\sigma_x)\}. \]
(ii)$_r$\ 对任意 $\sigma \in \{\sigma\}_r$，存在唯一的 $\pi \in \{\pi\}_r$，使得 $\sigma = \sigma_\pi$。
\end{theorem*}

\par\medskip\noindent\textbf{化归：}\
\begin{itemize}
\item $r = 1$：类域论；
\item (i) 中的唯一性：Chebotarev 密度定理；
\item \origpage{4}(ii) 中的唯一性：Piatetski-Shapiro 的强重数一定理。
\end{itemize}

% ------------------------------------------------------------
%  PAGE 5（原书 p. 4）
% ------------------------------------------------------------
\par\medskip\noindent\textbf{证明}\ \emph{对 $r$ 归纳}。

设 $r \ge 2$，且对 $r' < r$ 已知 (i)$_{r'}$、(ii)$_{r'}$。于是 (ii)$_r$ 也已知。这用到：
\begin{itemize}
\item Piatetski-Shapiro 的“互反定理”（配对 $L$-函数的论证）；
\item $r' < r$ 时的断言 (i)$_{r'}$；
\item Grothendieck 函数方程；
\item Laumon 乘积公式。
\end{itemize}

剩下要构造
\[ \pi \mapsto \sigma_\pi, \qquad \{\pi\}_r \to \{\sigma\}_r. \]

$\{\pi\}_r$ 中的 $\pi$ 也是 Hecke 代数上的不可约模，其中
\[ \mathcal{H}^r = C_c^\infty(\mathrm{GL}_r(\mathbb{A})) + \text{卷积} \]
（$\mathrm{dg}$ 为 Haar 测度，由 $\mathrm{dg}(K) = 1$ 归一化），
\[ K = K_\emptyset = \mathrm{GL}_r(O_{\mathbb{A}}) = \varprojlim_{N} \mathrm{GL}_r(\mathcal{O}_N), \]
\[ N \underset{\text{有限}}{=} \mathrm{Spec}\,\mathcal{O}_N \underset{\text{闭}}{\hookrightarrow} X \ :\ \text{水平}. \]

对每个水平 $N$，记
\[ K_N = \mathrm{Ker}[K \to \mathrm{GL}_r(\mathcal{O}_N)]. \]

$\mathcal{H}^r$ 是含幺子代数 $\mathcal{H}^r_N$ 的递增并，其中
\[ \mathcal{H}^r_N = C_c(K_N\backslash \mathrm{GL}_r(\mathbb{A})/K_N) + \text{卷积}. \]

记
\[ \{\pi\}^N_r = \{\pi \in \{\pi\}_r \mid \pi^{K_N} \neq 0 \quad \text{且} \quad \chi_\pi(a) = 1\}. \]

这是有限集。

% ------------------------------------------------------------
%  PAGE 6（原书 p. 5）
% ------------------------------------------------------------
\origpage{5}若 $\pi \in \{\pi\}^N_r$，则 $\pi^{K_N}$ 是 $\mathcal{H}^r_N$ 上的不可约模，且它刻画 $\pi$。只需构造
\[ \pi \mapsto \sigma_\pi, \qquad \{\pi\}^N_r \to \{\sigma\}_r. \]

为此，我们考察秩 $r$、带水平结构 $N$ 的 Drinfeld chtouca 之叠
\[ \mathrm{Cht}^r_N. \]
它们带有
\begin{itemize}
\item 一个维数为 $2r - 2$ 的光滑态射
\[ (\infty, 0) : \mathrm{Cht}^r_N \to (X - N) \times (X - N), \]
\item 群
\[ F^\times\backslash \mathbb{A}^\times \qquad \text{与} \qquad \mathrm{GL}_r(O_N) \]
的作用，该作用延拓为 Hecke 代数
\[ \mathcal{H}^r_N \]
经对应给出的作用。
\end{itemize}

\par\medskip\noindent\textbf{结论：}\ 若 $F^2 = F(X \times X)$ 为 $X \times X$ 上的函数域，则位于 $X \times X$ 泛点之上的紧支集 $\ell$-进上同调空间
\[ H^*_c(\mathrm{Cht}^r_N) \]
带有 $\mathcal{H}^r_N$ 与 $G_{F^2}$ 的作用。

\par\medskip\noindent\textbf{我们的目标：}\ 设 $q', q'' : X \times X \rightrightarrows X$ 为两个投影，在
\[ H^{2r-2}_c(\mathrm{Cht}^r_N/a^{\mathbb{Z}}) \]
中分离出如下形式的一块：
\[ \bigoplus_{\pi \in \{\pi\}^N_r} \pi^{K_N} \otimes q'^*\sigma_\pi \otimes q''^*\tilde{\sigma}_\pi(1-r). \]

\par\medskip\noindent\textbf{$\mathrm{Cht}^r$ 与 $\mathrm{Cht}^r_N$ 的定义：}\ 它们是相应模问题的解。

\begin{definition*}[Drinfeld]
(i)\ 对任意 $\mathbb{F}_q$-概形 $S$，$S$ 上秩 $r$ 的 chtouca $\tilde{\mathcal{E}}$ 由下列数据组成：
\begin{itemize}
\item \origpage{6}一个 $X \times S$ 上秩 $r$ 的丛 $\mathcal{E}$；
\item $\mathcal{E}$ 的一个修正，即一个交换图
\[ \mathcal{E} \xrightarrow{j} \mathcal{E}' \xleftarrow{t} \mathcal{E}'' \]
其中 $\mathcal{E}', \mathcal{E}''$ 是 $X \times S$ 上秩 $r$ 的丛，$j, t$ 是嵌入，使得 $\mathrm{Coker}\, j$ 与 $\mathrm{Coker}\, t$ 支撑于两个态射
\[ \infty, 0 : S \to X \]
的图像之上，并且在 $O_S$ 上可逆；
\item 一个同构
\[ \tau\mathcal{E} = (\mathrm{Id}_X \times \mathrm{Frob}_S)^*\mathcal{E} \xrightarrow{\sim} \mathcal{E}''. \]
\end{itemize}

(ii)\ 若 $\infty, 0 : S \to X$ 避开 $N$，则 $\mathcal{E}$ 上的水平结构 $N$ 是一个同构
\[ v : \mathcal{E}_N = \mathcal{E} \otimes_{O_{X\times S}} O_{N\times S} = \mathcal{E} \otimes_{O_X} O_N \xrightarrow{\sim} O^r_{N\times S} \]
使下图交换：
% 原文此处为交换图：上排 $\mathcal{E}_N \xrightarrow{\sim} \mathcal{E}'_N \xleftarrow{\sim} \mathcal{E}''_N \xleftarrow{\sim} \tau\mathcal{E}_N$，
% 两条斜箭头分别标 $v$、$\tau v$，同汇于 $O^r_{N\times S}$。
\begin{center}
\begin{tikzcd}[column sep=3.4em, row sep=2.6em]
\mathcal{E}_N \arrow[r,"\sim"] \arrow[drr,"v"'] & \mathcal{E}'_N
  & \mathcal{E}''_N \arrow[l,"\sim"']
  & \tau\mathcal{E}_N \arrow[l,"\sim"'] \arrow[dl,"\tau v"] \\
& & O^r_{N\times S}
\end{tikzcd}
\end{center}
\end{definition*}

\[
\begin{aligned}
\mathrm{Cht}^r &= \text{秩 $r$ 的 chtouca 之叠}\\
&= \text{Deligne--Mumford 意义下的代数叠}\\
(\infty, 0) &= \mathrm{Cht}^r \to X \times X \text{，是维数为 } 2r - 2 \text{ 的光滑态射}\\
\mathrm{Cht}^r_N &= \text{秩 $r$、带水平结构 $N$ 的 chtouca 之叠}\\
&= \text{Deligne--Mumford 意义下的代数叠。}
\end{aligned}
\]

% ------------------------------------------------------------
%  PAGE 8（原书 p. 7）
% ------------------------------------------------------------
\origpage{7}
$\mathrm{Cht}^r_N \to \mathrm{Cht}^r \times_{X\times X} (X - N) \times (X - N)$ 是可表的、有限、étale Galois 的态射，其群为 $\mathrm{GL}_r(O_N)$。

\par\medskip\noindent\textbf{$F^\times\backslash \mathbb{A}^\times$ 的作用：}\
\[ F^\times\backslash \mathbb{A}^\times / O^\times_{\mathbb{A}} \cong \{\text{$X$ 上的可逆层 } \mathcal{L}\} \]
\[ F^\times\backslash \mathbb{A}^\times / (O^\times_{\mathbb{A}} \cap K_N) \cong \{\text{$X$ 上的可逆层 } \mathcal{L} \text{，带水平结构 } N\} \]
由此得到在 $\tilde{\mathcal{E}} = (\mathcal{E} \hookrightarrow \mathcal{E}' \hookleftarrow \mathcal{E}'' \xleftarrow{\sim} \tau\mathcal{E})$ 上的作用：
\[ (\mathcal{L}, \tilde{\mathcal{E}}) \mapsto \mathcal{L} \otimes \tilde{\mathcal{E}}. \]

\par\medskip\noindent\textbf{$\mathcal{H}^r_\emptyset$ 在 $\mathrm{Cht}^r$ 上的作用：}\
只需明确与双陪集 $K \cdot g \cdot K$ 相伴的对应，其中 $g = (g_x)_{x \in |X|} \in \mathrm{GL}_r(\mathbb{A})$。

设 $T \subset |X|$ 是有限子集，使得
\[ x \notin T \Rightarrow g_x \in \mathrm{GL}_r(O_x). \]

设 $\tilde{\mathcal{E}}$ 是 $\overline{\mathbb{F}}_q$ 上的 chtouca，其零点与极点都在 $U = X - T$ 中。

对 $x \in T$，$\mathcal{E} \otimes_{O_X} O_x$ 定义了 $\mathcal{E} \otimes_{O_X} F_x$ 中一个在 Frob 下稳定的格；并且存在保持 Frob 的同构
\[ \mathcal{E} \otimes_{O_X} F_x \cong (F_x \otimes_{\mathbb{F}_q} \overline{\mathbb{F}}_q)^r \]
使得 $\mathcal{E} \otimes_{O_X} O_x$ 来自 $F^r_x$ 的一个格。

设 $\tilde{\mathcal{F}}$ 是另一个 chtouca，使得 $\mathcal{E} \otimes_{O_X} F_x$ 与 $\mathcal{F} \otimes_{O_X} F_x$ 等同。两个格
\[ \mathcal{E} \otimes_{O_X} O_x \qquad \text{与} \qquad \mathcal{F} \otimes_{O_X} O_x \]
定义了
\[ \mathrm{GL}_r(O_x)\backslash \mathrm{GL}_r(F_x)/\mathrm{GL}_r(O_x) \]
中的一个元素，它不依赖于同构
\[ \mathcal{E} \otimes_{O_X} F_x = \mathcal{F} \otimes_{O_X} F_x \cong (F_x \otimes_{\mathbb{F}_q} \overline{\mathbb{F}}_q)^r. \]

对应 $K \cdot g \cdot K$ 的一个点 $(\tilde{\mathcal{E}}, \tilde{\mathcal{F}})$ 由下列数据组成：

% ------------------------------------------------------------
%  PAGE 9（原书 p. 8）
% ------------------------------------------------------------
\origpage{8}
\begin{itemize}
\item 两个 chtouca $\tilde{\mathcal{E}}$、$\tilde{\mathcal{F}}$，其零点与极点都在 $X - T$ 中；
\item 它们限制在 $X - T = U$ 上的一个等同（由此诱导出等同 $\mathcal{E} \otimes_{O_X} F_x = \mathcal{F} \otimes_{O_X} F_x$，$\forall x$），
\end{itemize}
使得对每个 $x \in T$，格对
\[ (\mathcal{E} \otimes_{O_X} O_x, \mathcal{F} \otimes_{O_X} O_x) \]
都定义出元素 $g_x \in \mathrm{GL}_r(O_x)\backslash \mathrm{GL}_r(F_x)/\mathrm{GL}_r(O_x)$。

\par\medskip\noindent\textbf{非常重要的注：}\ 若 $\tilde{\mathcal{E}}$ 与 $\tilde{\mathcal{F}}$ 是 Hecke 意义下互相对应的两个 chtouca，则它们在 $X$ 的一个非空开集上等同。

注意到
\[ \mathrm{Cht}^r = \coprod_{d \in \mathbb{Z}} \mathrm{Cht}^{r,d} \]
其中每个 $\mathrm{Cht}^{r,d}$ 只有有限多个连通分支。因此，我们作
\[ \mathrm{Cht}^r/a^{\mathbb{Z}} \cong \coprod_{1 \le d \le r|\deg(a)|} \mathrm{Cht}^{r,d}. \]

它有有限多个连通分支，并且是局部有限型的，但
\[ \emph{它不是有限型的} \]
$\rightsquigarrow$ $H^*_c(\mathrm{Cht}^r_N/a^{\mathbb{Z}})$ 是无限维的；

$\rightsquigarrow$ 对 $x \in |(X - N) \times (X - N)|$、$\deg(x) \mid s$ 与 $f \in \mathcal{H}^r_N$，不动点个数
\[ \mathrm{Lef}_x(f \times \mathrm{Frob}^s, \mathrm{Cht}^r_N/a^{\mathbb{Z}}) \]
是无限的。

\par\medskip\noindent\textbf{用 Harder--Narasimhan 典范多边形作截断}\

若 $\tilde{\mathcal{E}} = (\mathcal{E} \hookrightarrow \mathcal{E}' \hookleftarrow \tau\mathcal{E})$ 是 $\overline{\mathbb{F}}_q$ 上的 chtouca，则 $\tilde{\mathcal{E}}$ 的子对象 $\tilde{\mathcal{F}}$ 由两个同秩的子丛 $\mathcal{F} \hookrightarrow \mathcal{E}$、$\mathcal{F}' \hookrightarrow \mathcal{E}'$ 组成，使得 $\mathcal{F} \hookrightarrow \mathcal{F}'$ 且 $\tau\mathcal{F} \hookrightarrow \mathcal{F}'$。

可以令
\[ \mathrm{rg}\,\tilde{\mathcal{F}} = \mathrm{rg}\,\mathcal{F} = \mathrm{rg}\,\mathcal{F}' \]
并令 $\deg\tilde{\mathcal{F}} = \deg\mathcal{F}$。

% ------------------------------------------------------------
%  PAGE 10（原书 p. 9）
% ------------------------------------------------------------
\origpage{9}
% 原文此处有图：横轴标 $0$ 与 $r$，折线自原点上升，经三个标出的顶点，过最高点后下降回到 $r$。
\begin{center}
\begin{tikzpicture}[x=1.05cm, y=1.05cm]
  \draw (-0.25,0) -- (5.5,0);
  \draw (0,-0.25) -- (0,2.05);
  \draw (0,0) -- (1.15,0.80) -- (2.45,1.30) -- (3.35,0.92) -- (4.70,0);
  \fill (1.15,0.80) circle (1.7pt);
  \fill (2.45,1.30) circle (1.7pt);
  \fill (3.35,0.92) circle (1.7pt);
  \node[below left] at (0,0) {$0$};
  \node[below] at (4.70,0) {$r$};
\end{tikzpicture}
\end{center}

所谓多边形，是指满足下述条件的映射 $p : [0, r] \to \mathbb{R}$：
\begin{itemize}
\item $p(0) = p(r) = 0$；
\item $p$ 在每个区间 $[r' - 1, r']$ 上是仿射的，$0 < r' \le r$。
\end{itemize}

$\tilde{\mathcal{E}}$ 的 Harder--Narasimhan 典范多边形 $\overline{p}$ 是满足下述条件的最小多边形：对 $\tilde{\mathcal{E}}$ 的每个子对象 $\tilde{\mathcal{F}}$，都有
\[ \deg\tilde{\mathcal{F}} - \frac{\mathrm{rg}\,\tilde{\mathcal{F}}}{r}\deg\tilde{\mathcal{E}} \le p(\mathrm{rg}\,\tilde{\mathcal{F}}). \]
它是凸的。

\begin{proposition*}
设 $p$ 是凸的截断多边形，则
\[
\begin{aligned}
\mathrm{Cht}^{r,d,\overline{p}\le p} &\subset \mathrm{Cht}^{r,d} & \mathrm{Cht}^{r,d,\overline{p}\le p}_N\\
\mathrm{Cht}^{r,\overline{p}\le p}/a^{\mathbb{Z}} &\subset \mathrm{Cht}^r/a^{\mathbb{Z}} & \mathrm{Cht}^{r,\overline{p}\le p}_N/a^{\mathbb{Z}}
\end{aligned}
\]
都是有限型的开子叠。\footnote{译注：原文该展示右列两式之前未印 $\subset$ 符号，此处照录。}
\end{proposition*}

\begin{theorem*}
设 $x \in |(X - N) \times (X - N)|$ 映到 $\infty \neq 0$，$s \in \mathbb{N}$ 满足 $\deg(x) \mid s$，$f \in \mathcal{H}^r_N$，则有不动点计数公式
\[
\begin{aligned}
&\frac{1}{r!}\sum_{k=1}^{r!} \mathrm{Lef}_{(\mathrm{Frob}^k_X \times \mathrm{Id}_X)(x)}\big(f \times \mathrm{Frob}^s, \mathrm{Cht}^{r,\overline{p}\le p}_N/a^{\mathbb{Z}}\big)\\[4pt]
&= q^{(r-1)s} \sum_{\pi \in \{\pi\}^N_r} \mathrm{Tr}_\pi(f)\big(z_1(\pi_\infty)^{-s/\deg(\infty)} + \cdots + z_r(\pi_\infty)^{-s/\deg(\infty)}\big)\\[4pt]
&\qquad\qquad \times \big(z_1(\pi_0)^{s/\deg(0)} + \cdots + z_r(\pi_0)^{s/\deg(0)}\big)
\end{aligned}
\]
+ 其他项，其中出现 $\mathrm{GL}_{r_1} \times \cdots \times \mathrm{GL}_{r_k}$ 的尖点自守表示的 Hecke 特征值，
\[ r_1 + \cdots + r_k = r \qquad (k \ge 2). \]
\end{theorem*}

% ============================================================
%  译者附记（不进入排版）
%
%  【\cite key】本段（原书 pp. 3--9）原文无方括号引用，故未使用 \cite。
%
%  【待定术语】术语表未收，暂按下述译法处理：
%    Conséquence      → 结论（“推论”已固定对应 corollaire，故避开）
%    réunion filtrante → 递增并
%    sous-algèbre unitaire → 含幺子代数
%    ouvert de type fini → 有限型开子叠
%    double classe   → 双陪集
%    identification  → 等同
%    réseau          → 格（随术语表 $\varphi$-réseau = $\varphi$-格）
%    Coker           → 保留 \mathrm{Coker}
%    Lef             → 保留 \mathrm{Lef}（Lefschetz 数）
%
%  【体例说明】
%    - 公式源码一律照录转录稿（含 $O_x$ 与 $\mathcal{O}_N$ 的原有区分，
%      原书即如此），仅把 \text{}/\emph{} 中的法文散文译为中文。
%    - 转录稿中“⟨diagramme p. 6⟩”一类描述性按语，已于相应处用 tikz-cd
%      或 tikz 重绘，不另出注。
% ============================================================
